% Options for packages loaded elsewhere
\PassOptionsToPackage{unicode}{hyperref}
\PassOptionsToPackage{hyphens}{url}
\documentclass[
  12pt,
]{article}
\usepackage{xcolor}
\usepackage[margin=2.5cm]{geometry}
\usepackage{amsmath,amssymb}
\setcounter{secnumdepth}{5}
\usepackage{iftex}
\ifPDFTeX
  \usepackage[T1]{fontenc}
  \usepackage[utf8]{inputenc}
  \usepackage{textcomp} % provide euro and other symbols
\else % if luatex or xetex
  \usepackage{unicode-math} % this also loads fontspec
  \defaultfontfeatures{Scale=MatchLowercase}
  \defaultfontfeatures[\rmfamily]{Ligatures=TeX,Scale=1}
\fi
\usepackage{lmodern}
\ifPDFTeX\else
  % xetex/luatex font selection
\fi
% Use upquote if available, for straight quotes in verbatim environments
\IfFileExists{upquote.sty}{\usepackage{upquote}}{}
\IfFileExists{microtype.sty}{% use microtype if available
  \usepackage[]{microtype}
  \UseMicrotypeSet[protrusion]{basicmath} % disable protrusion for tt fonts
}{}
\makeatletter
\@ifundefined{KOMAClassName}{% if non-KOMA class
  \IfFileExists{parskip.sty}{%
    \usepackage{parskip}
  }{% else
    \setlength{\parindent}{0pt}
    \setlength{\parskip}{6pt plus 2pt minus 1pt}}
}{% if KOMA class
  \KOMAoptions{parskip=half}}
\makeatother
\usepackage{color}
\usepackage{fancyvrb}
\newcommand{\VerbBar}{|}
\newcommand{\VERB}{\Verb[commandchars=\\\{\}]}
\DefineVerbatimEnvironment{Highlighting}{Verbatim}{commandchars=\\\{\}}
% Add ',fontsize=\small' for more characters per line
\usepackage{framed}
\definecolor{shadecolor}{RGB}{248,248,248}
\newenvironment{Shaded}{\begin{snugshade}}{\end{snugshade}}
\newcommand{\AlertTok}[1]{\textcolor[rgb]{0.94,0.16,0.16}{#1}}
\newcommand{\AnnotationTok}[1]{\textcolor[rgb]{0.56,0.35,0.01}{\textbf{\textit{#1}}}}
\newcommand{\AttributeTok}[1]{\textcolor[rgb]{0.13,0.29,0.53}{#1}}
\newcommand{\BaseNTok}[1]{\textcolor[rgb]{0.00,0.00,0.81}{#1}}
\newcommand{\BuiltInTok}[1]{#1}
\newcommand{\CharTok}[1]{\textcolor[rgb]{0.31,0.60,0.02}{#1}}
\newcommand{\CommentTok}[1]{\textcolor[rgb]{0.56,0.35,0.01}{\textit{#1}}}
\newcommand{\CommentVarTok}[1]{\textcolor[rgb]{0.56,0.35,0.01}{\textbf{\textit{#1}}}}
\newcommand{\ConstantTok}[1]{\textcolor[rgb]{0.56,0.35,0.01}{#1}}
\newcommand{\ControlFlowTok}[1]{\textcolor[rgb]{0.13,0.29,0.53}{\textbf{#1}}}
\newcommand{\DataTypeTok}[1]{\textcolor[rgb]{0.13,0.29,0.53}{#1}}
\newcommand{\DecValTok}[1]{\textcolor[rgb]{0.00,0.00,0.81}{#1}}
\newcommand{\DocumentationTok}[1]{\textcolor[rgb]{0.56,0.35,0.01}{\textbf{\textit{#1}}}}
\newcommand{\ErrorTok}[1]{\textcolor[rgb]{0.64,0.00,0.00}{\textbf{#1}}}
\newcommand{\ExtensionTok}[1]{#1}
\newcommand{\FloatTok}[1]{\textcolor[rgb]{0.00,0.00,0.81}{#1}}
\newcommand{\FunctionTok}[1]{\textcolor[rgb]{0.13,0.29,0.53}{\textbf{#1}}}
\newcommand{\ImportTok}[1]{#1}
\newcommand{\InformationTok}[1]{\textcolor[rgb]{0.56,0.35,0.01}{\textbf{\textit{#1}}}}
\newcommand{\KeywordTok}[1]{\textcolor[rgb]{0.13,0.29,0.53}{\textbf{#1}}}
\newcommand{\NormalTok}[1]{#1}
\newcommand{\OperatorTok}[1]{\textcolor[rgb]{0.81,0.36,0.00}{\textbf{#1}}}
\newcommand{\OtherTok}[1]{\textcolor[rgb]{0.56,0.35,0.01}{#1}}
\newcommand{\PreprocessorTok}[1]{\textcolor[rgb]{0.56,0.35,0.01}{\textit{#1}}}
\newcommand{\RegionMarkerTok}[1]{#1}
\newcommand{\SpecialCharTok}[1]{\textcolor[rgb]{0.81,0.36,0.00}{\textbf{#1}}}
\newcommand{\SpecialStringTok}[1]{\textcolor[rgb]{0.31,0.60,0.02}{#1}}
\newcommand{\StringTok}[1]{\textcolor[rgb]{0.31,0.60,0.02}{#1}}
\newcommand{\VariableTok}[1]{\textcolor[rgb]{0.00,0.00,0.00}{#1}}
\newcommand{\VerbatimStringTok}[1]{\textcolor[rgb]{0.31,0.60,0.02}{#1}}
\newcommand{\WarningTok}[1]{\textcolor[rgb]{0.56,0.35,0.01}{\textbf{\textit{#1}}}}
\usepackage{longtable,booktabs,array}
\usepackage{calc} % for calculating minipage widths
% Correct order of tables after \paragraph or \subparagraph
\usepackage{etoolbox}
\makeatletter
\patchcmd\longtable{\par}{\if@noskipsec\mbox{}\fi\par}{}{}
\makeatother
% Allow footnotes in longtable head/foot
\IfFileExists{footnotehyper.sty}{\usepackage{footnotehyper}}{\usepackage{footnote}}
\makesavenoteenv{longtable}
\usepackage{graphicx}
\makeatletter
\newsavebox\pandoc@box
\newcommand*\pandocbounded[1]{% scales image to fit in text height/width
  \sbox\pandoc@box{#1}%
  \Gscale@div\@tempa{\textheight}{\dimexpr\ht\pandoc@box+\dp\pandoc@box\relax}%
  \Gscale@div\@tempb{\linewidth}{\wd\pandoc@box}%
  \ifdim\@tempb\p@<\@tempa\p@\let\@tempa\@tempb\fi% select the smaller of both
  \ifdim\@tempa\p@<\p@\scalebox{\@tempa}{\usebox\pandoc@box}%
  \else\usebox{\pandoc@box}%
  \fi%
}
% Set default figure placement to htbp
\def\fps@figure{htbp}
\makeatother
\ifLuaTeX
\usepackage[bidi=basic]{babel}
\else
\usepackage[bidi=default]{babel}
\fi
% get rid of language-specific shorthands (see #6817):
\let\LanguageShortHands\languageshorthands
\def\languageshorthands#1{}
\setlength{\emergencystretch}{3em} % prevent overfull lines
\providecommand{\tightlist}{%
  \setlength{\itemsep}{0pt}\setlength{\parskip}{0pt}}
\\usepackage[english]{babel}
\\usepackage{amsmath,amssymb,amsthm}
\\usepackage{booktabs}
\\usepackage{float}
\usepackage{bookmark}
\IfFileExists{xurl.sty}{\usepackage{xurl}}{} % add URL line breaks if available
\urlstyle{same}
\hypersetup{
  pdftitle={Bayesian Analysis Task III},
  pdfauthor={Leonel Fernando Nabaza Ruibal},
  pdflang={eng},
  hidelinks,
  pdfcreator={LaTeX via pandoc}}

\title{Bayesian Analysis Task III}
\author{Leonel Fernando Nabaza Ruibal}
\date{11/04/2026}

\begin{document}
\maketitle

{
\setcounter{tocdepth}{2}
\tableofcontents
}
\section{Wording and data}\label{wording-and-data}

We want to study whether the final grades in the \emph{Bayesian
Analysis} course are the same for boys and girls.

The observed data are:

\begin{Shaded}
\begin{Highlighting}[]
\NormalTok{boys  }\OtherTok{\textless{}{-}} \FunctionTok{c}\NormalTok{(}\FloatTok{9.6}\NormalTok{, }\FloatTok{7.0}\NormalTok{, }\FloatTok{5.0}\NormalTok{, }\FloatTok{8.0}\NormalTok{, }\FloatTok{8.4}\NormalTok{, }\FloatTok{6.4}\NormalTok{)}
\NormalTok{girls }\OtherTok{\textless{}{-}} \FunctionTok{c}\NormalTok{(}\FloatTok{6.1}\NormalTok{, }\FloatTok{9.1}\NormalTok{, }\FloatTok{8.8}\NormalTok{, }\FloatTok{5.7}\NormalTok{, }\FloatTok{8.9}\NormalTok{, }\FloatTok{6.1}\NormalTok{, }\FloatTok{6.5}\NormalTok{)}

\NormalTok{df }\OtherTok{\textless{}{-}} \FunctionTok{data.frame}\NormalTok{(}
  \AttributeTok{grade =} \FunctionTok{c}\NormalTok{(boys, girls),}
  \AttributeTok{sex   =} \FunctionTok{c}\NormalTok{(}\FunctionTok{rep}\NormalTok{(}\DecValTok{1}\NormalTok{, }\FunctionTok{length}\NormalTok{(boys)), }\FunctionTok{rep}\NormalTok{(}\DecValTok{0}\NormalTok{, }\FunctionTok{length}\NormalTok{(girls))),}
  \AttributeTok{group =} \FunctionTok{c}\NormalTok{(}\FunctionTok{rep}\NormalTok{(}\StringTok{"boys"}\NormalTok{, }\FunctionTok{length}\NormalTok{(boys)), }\FunctionTok{rep}\NormalTok{(}\StringTok{"girls"}\NormalTok{, }\FunctionTok{length}\NormalTok{(girls)))}
\NormalTok{)}

\NormalTok{knitr}\SpecialCharTok{::}\FunctionTok{kable}\NormalTok{(df, }\AttributeTok{col.names =} \FunctionTok{c}\NormalTok{(}\StringTok{"grade"}\NormalTok{, }\StringTok{"sex"}\NormalTok{, }\StringTok{"group"}\NormalTok{), }\AttributeTok{booktabs =} \ConstantTok{TRUE}\NormalTok{)}
\end{Highlighting}
\end{Shaded}

\begin{longtable}[]{@{}rrl@{}}
\toprule\noalign{}
grade & sex & group \\
\midrule\noalign{}
\endhead
\bottomrule\noalign{}
\endlastfoot
9.6 & 1 & boys \\
7.0 & 1 & boys \\
5.0 & 1 & boys \\
8.0 & 1 & boys \\
8.4 & 1 & boys \\
6.4 & 1 & boys \\
6.1 & 0 & girls \\
9.1 & 0 & girls \\
8.8 & 0 & girls \\
5.7 & 0 & girls \\
8.9 & 0 & girls \\
6.1 & 0 & girls \\
6.5 & 0 & girls \\
\end{longtable}

The statement proposes the model \[
y_i \mid \beta_0, \beta_1, \sigma \sim \mathcal N(\beta_0 + \beta_1\, \text{sex}_i, \sigma),
\] where \(\text{sex}_i = 1\) for boys and \(\text{sex}_i = 0\) for
girls.

\section{Initial exploration}\label{initial-exploration}

\begin{Shaded}
\begin{Highlighting}[]
\NormalTok{summary\_by\_group }\OtherTok{\textless{}{-}} \FunctionTok{aggregate}\NormalTok{(grade }\SpecialCharTok{\textasciitilde{}}\NormalTok{ group, }\AttributeTok{data =}\NormalTok{ df, }\ControlFlowTok{function}\NormalTok{(x) }\FunctionTok{c}\NormalTok{(}\AttributeTok{n =} \FunctionTok{length}\NormalTok{(x), }\AttributeTok{mean =} \FunctionTok{mean}\NormalTok{(x), }\AttributeTok{sd =} \FunctionTok{sd}\NormalTok{(x)))}
\NormalTok{summary\_by\_group}
\end{Highlighting}
\end{Shaded}

\begin{verbatim}
##   group  grade.n grade.mean grade.sd
## 1  boys 6.000000   7.400000 1.619877
## 2 girls 7.000000   7.314286 1.534523
\end{verbatim}

\begin{Shaded}
\begin{Highlighting}[]
\FunctionTok{library}\NormalTok{(ggplot2)}

\FunctionTok{ggplot}\NormalTok{(df, }\FunctionTok{aes}\NormalTok{(}\AttributeTok{x =}\NormalTok{ group, }\AttributeTok{y =}\NormalTok{ grade)) }\SpecialCharTok{+}
  \FunctionTok{geom\_boxplot}\NormalTok{(}\AttributeTok{outlier.shape =} \ConstantTok{NA}\NormalTok{) }\SpecialCharTok{+}
  \FunctionTok{geom\_jitter}\NormalTok{(}\AttributeTok{width =} \FloatTok{0.1}\NormalTok{, }\AttributeTok{height =} \DecValTok{0}\NormalTok{) }\SpecialCharTok{+}
  \FunctionTok{labs}\NormalTok{(}\AttributeTok{title =} \StringTok{"Grades by group"}\NormalTok{, }\AttributeTok{x =} \StringTok{"Group"}\NormalTok{, }\AttributeTok{y =} \StringTok{"Grade"}\NormalTok{)}
\end{Highlighting}
\end{Shaded}

\begin{center}\includegraphics{bayesian_grades_stan_files/figure-latex/data-boxplot-1} \end{center}

At first glance, the sample means are very similar, but the dispersion
does not appear to be exactly the same in both groups. Moreover, grades
are bounded between 0 and 10, so a pure normal model does not strictly
respect the support of the variable.

\section{a) Full Bayesian model}\label{a-full-bayesian-model}

\subsection{Model proposed in the
statement}\label{model-proposed-in-the-statement}

The linear model in the statement can be written as \[
y_i \mid \beta_0, \beta_1, \sigma \overset{ind}{\sim} \mathcal N(\beta_0 + \beta_1 x_i, \sigma), \qquad x_i \in \{0,1\}.
\]

Here:

\begin{itemize}
\tightlist
\item
  \(\beta_0\) is the mean grade for girls.
\item
  \(\beta_0 + \beta_1\) is the mean grade for boys.
\item
  \(\beta_1\) is the difference in means, boys minus girls.
\item
  \(\sigma\) is the common standard deviation for both groups.
\end{itemize}

A reasonable choice of weakly informative priors is \[
\beta_0 \sim \mathcal N(7, 2^2), \qquad
\beta_1 \sim \mathcal N(0, 2^2), \qquad
\sigma \sim \text{Half-Normal}(0, 2).
\]

\subsubsection{Justification of the
priors}\label{justification-of-the-priors}

\begin{itemize}
\tightlist
\item
  Grades are on a 0--10 scale, so centering \(\beta_0\) around 7 is
  plausible and still weakly informative.
\item
  For \(\beta_1\), centering at 0 expresses that \emph{a priori} we do
  not favor either group having better grades.
\item
  The half-normal prior for \(\sigma\) guarantees positivity and avoids
  excessively heavy tails in such a small sample.
\end{itemize}

\subsection{Prior plots}\label{prior-plots}

\begin{Shaded}
\begin{Highlighting}[]
\FunctionTok{par}\NormalTok{(}\AttributeTok{mfrow =} \FunctionTok{c}\NormalTok{(}\DecValTok{1}\NormalTok{, }\DecValTok{3}\NormalTok{))}
\FunctionTok{curve}\NormalTok{(}\FunctionTok{dnorm}\NormalTok{(x, }\AttributeTok{mean =} \DecValTok{7}\NormalTok{, }\AttributeTok{sd =} \DecValTok{2}\NormalTok{), }\AttributeTok{from =} \SpecialCharTok{{-}}\DecValTok{1}\NormalTok{, }\AttributeTok{to =} \DecValTok{12}\NormalTok{,}
      \AttributeTok{main =} \FunctionTok{expression}\NormalTok{(}\FunctionTok{paste}\NormalTok{(}\StringTok{"Prior for "}\NormalTok{, beta[}\DecValTok{0}\NormalTok{])), }\AttributeTok{xlab =} \FunctionTok{expression}\NormalTok{(beta[}\DecValTok{0}\NormalTok{]), }\AttributeTok{ylab =} \StringTok{"density"}\NormalTok{)}
\FunctionTok{curve}\NormalTok{(}\FunctionTok{dnorm}\NormalTok{(x, }\AttributeTok{mean =} \DecValTok{0}\NormalTok{, }\AttributeTok{sd =} \DecValTok{2}\NormalTok{), }\AttributeTok{from =} \SpecialCharTok{{-}}\DecValTok{6}\NormalTok{, }\AttributeTok{to =} \DecValTok{6}\NormalTok{,}
      \AttributeTok{main =} \FunctionTok{expression}\NormalTok{(}\FunctionTok{paste}\NormalTok{(}\StringTok{"Prior for "}\NormalTok{, beta[}\DecValTok{1}\NormalTok{])), }\AttributeTok{xlab =} \FunctionTok{expression}\NormalTok{(beta[}\DecValTok{1}\NormalTok{]), }\AttributeTok{ylab =} \StringTok{"density"}\NormalTok{)}
\NormalTok{xs }\OtherTok{\textless{}{-}} \FunctionTok{seq}\NormalTok{(}\DecValTok{0}\NormalTok{, }\DecValTok{6}\NormalTok{, }\AttributeTok{length.out =} \DecValTok{500}\NormalTok{)}
\FunctionTok{plot}\NormalTok{(xs, }\DecValTok{2}\SpecialCharTok{*}\FunctionTok{dnorm}\NormalTok{(xs, }\AttributeTok{mean =} \DecValTok{0}\NormalTok{, }\AttributeTok{sd =} \DecValTok{2}\NormalTok{), }\AttributeTok{type =} \StringTok{"l"}\NormalTok{,}
     \AttributeTok{main =} \FunctionTok{expression}\NormalTok{(}\FunctionTok{paste}\NormalTok{(}\StringTok{"Prior for "}\NormalTok{, sigma)), }\AttributeTok{xlab =} \FunctionTok{expression}\NormalTok{(sigma), }\AttributeTok{ylab =} \StringTok{"density"}\NormalTok{)}
\end{Highlighting}
\end{Shaded}

\begin{center}\includegraphics{bayesian_grades_stan_files/figure-latex/prior-plots-statement-model-1} \end{center}

\begin{Shaded}
\begin{Highlighting}[]
\FunctionTok{par}\NormalTok{(}\AttributeTok{mfrow =} \FunctionTok{c}\NormalTok{(}\DecValTok{1}\NormalTok{, }\DecValTok{1}\NormalTok{))}
\end{Highlighting}
\end{Shaded}

\section{b) Discussion of the model and alternative
proposal}\label{b-discussion-of-the-model-and-alternative-proposal}

The model in the statement is useful as a starting point, but it has two
important limitations:

\begin{enumerate}
\def\labelenumi{\arabic{enumi}.}
\tightlist
\item
  \textbf{Common variance}. It imposes the same standard deviation for
  boys and girls. With so few observations, this restriction may be too
  strong.
\item
  \textbf{Unbounded support}. The normal distribution allows values
  below 0 and above 10, which does not perfectly match the nature of
  grades.
\end{enumerate}

\subsection{Proposed alternative
model}\label{proposed-alternative-model}

I propose using a model with \textbf{different means and different
standard deviations} for each group: \[
y_{g,j} \mid \mu_g, \sigma_g \sim \mathcal N(\mu_g, \sigma_g),
\qquad
y_{b,k} \mid \mu_b, \sigma_b \sim \mathcal N(\mu_b, \sigma_b).
\]

With weakly informative priors: \[
\mu_g \sim \mathcal N(7, 2^2), \qquad
\mu_b \sim \mathcal N(7, 2^2), \qquad
\sigma_g, \sigma_b \sim \text{Half-Normal}(0, 2).
\]

This model remains simple, but it is more flexible than the original
one. If we wanted to be even more realistic with the support \([0,10]\),
we could use a truncated normal model or a beta model after rescaling
the grades to \((0,1)\). However, for this sample size, the normal model
with separate variances is a very reasonable and easy option to
implement in Stan.

\section{Chosen model for Stan}\label{chosen-model-for-stan}

\subsection{Stan code}\label{stan-code}

Below we use the model with separate means and separate standard
deviations. In addition, in \texttt{generated\ quantities} we request:

\begin{itemize}
\tightlist
\item
  the difference in means \(\delta = \mu_b - \mu_g\),
\item
  a posterior predictive draw for a new girl, \texttt{y\_new\_girl}.
\end{itemize}

\begin{Shaded}
\begin{Highlighting}[]
\NormalTok{stan\_code }\OtherTok{\textless{}{-}} \StringTok{\textquotesingle{}}
\StringTok{data \{}
\StringTok{  int\textless{}lower=1\textgreater{} N\_boys;}
\StringTok{  int\textless{}lower=1\textgreater{} N\_girls;}
\StringTok{  vector[N\_boys] y\_boys;}
\StringTok{  vector[N\_girls] y\_girls;}
\StringTok{\}}
\StringTok{parameters \{}
\StringTok{  real mu\_boys;}
\StringTok{  real mu\_girls;}
\StringTok{  real\textless{}lower=0\textgreater{} sigma\_boys;}
\StringTok{  real\textless{}lower=0\textgreater{} sigma\_girls;}
\StringTok{\}}
\StringTok{model \{}
\StringTok{  // priors}
\StringTok{  mu\_boys \textasciitilde{} normal(7, 2);}
\StringTok{  mu\_girls \textasciitilde{} normal(7, 2);}
\StringTok{  sigma\_boys \textasciitilde{} normal(0, 2);}
\StringTok{  sigma\_girls \textasciitilde{} normal(0, 2);}

\StringTok{  // likelihood}
\StringTok{  y\_boys \textasciitilde{} normal(mu\_boys, sigma\_boys);}
\StringTok{  y\_girls \textasciitilde{} normal(mu\_girls, sigma\_girls);}
\StringTok{\}}
\StringTok{generated quantities \{}
\StringTok{  real delta;}
\StringTok{  real y\_new\_girl;}
\StringTok{  delta = mu\_boys {-} mu\_girls;}
\StringTok{  y\_new\_girl = normal\_rng(mu\_girls, sigma\_girls);}
\StringTok{\}}
\StringTok{\textquotesingle{}}
\FunctionTok{writeLines}\NormalTok{(stan\_code, }\StringTok{"bayesian\_grades\_model.stan"}\NormalTok{)}
\FunctionTok{cat}\NormalTok{(stan\_code)}
\end{Highlighting}
\end{Shaded}

\begin{verbatim}
## 
## data {
##   int<lower=1> N_boys;
##   int<lower=1> N_girls;
##   vector[N_boys] y_boys;
##   vector[N_girls] y_girls;
## }
## parameters {
##   real mu_boys;
##   real mu_girls;
##   real<lower=0> sigma_boys;
##   real<lower=0> sigma_girls;
## }
## model {
##   // priors
##   mu_boys ~ normal(7, 2);
##   mu_girls ~ normal(7, 2);
##   sigma_boys ~ normal(0, 2);
##   sigma_girls ~ normal(0, 2);
## 
##   // likelihood
##   y_boys ~ normal(mu_boys, sigma_boys);
##   y_girls ~ normal(mu_girls, sigma_girls);
## }
## generated quantities {
##   real delta;
##   real y_new_girl;
##   delta = mu_boys - mu_girls;
##   y_new_girl = normal_rng(mu_girls, sigma_girls);
## }
\end{verbatim}

\subsection{\texorpdfstring{Fitting with
\texttt{rstan}}{Fitting with rstan}}\label{fitting-with-rstan}

\begin{Shaded}
\begin{Highlighting}[]
\FunctionTok{library}\NormalTok{(rstan)}
\FunctionTok{library}\NormalTok{(bayesplot)}
\FunctionTok{library}\NormalTok{(posterior)}

\FunctionTok{rstan\_options}\NormalTok{(}\AttributeTok{auto\_write =} \ConstantTok{TRUE}\NormalTok{)}
\FunctionTok{options}\NormalTok{(}\AttributeTok{mc.cores =}\NormalTok{ parallel}\SpecialCharTok{::}\FunctionTok{detectCores}\NormalTok{())}

\NormalTok{stan\_data }\OtherTok{\textless{}{-}} \FunctionTok{list}\NormalTok{(}
  \AttributeTok{N\_boys  =} \FunctionTok{length}\NormalTok{(boys),}
  \AttributeTok{N\_girls =} \FunctionTok{length}\NormalTok{(girls),}
  \AttributeTok{y\_boys  =}\NormalTok{ boys,}
  \AttributeTok{y\_girls =}\NormalTok{ girls}
\NormalTok{)}

\NormalTok{fit }\OtherTok{\textless{}{-}} \FunctionTok{stan}\NormalTok{(}
  \AttributeTok{file =} \StringTok{"bayesian\_grades\_model.stan"}\NormalTok{,}
  \AttributeTok{data =}\NormalTok{ stan\_data,}
  \AttributeTok{chains =} \DecValTok{4}\NormalTok{,}
  \AttributeTok{iter =} \DecValTok{4000}\NormalTok{,}
  \AttributeTok{warmup =} \DecValTok{2000}\NormalTok{,}
  \AttributeTok{seed =} \DecValTok{1234}\NormalTok{,}
  \AttributeTok{refresh =} \DecValTok{0}
\NormalTok{)}

\FunctionTok{print}\NormalTok{(fit, }\AttributeTok{pars =} \FunctionTok{c}\NormalTok{(}\StringTok{"mu\_boys"}\NormalTok{, }\StringTok{"mu\_girls"}\NormalTok{, }\StringTok{"sigma\_boys"}\NormalTok{, }\StringTok{"sigma\_girls"}\NormalTok{, }\StringTok{"delta"}\NormalTok{, }\StringTok{"y\_new\_girl"}\NormalTok{),}
      \AttributeTok{probs =} \FunctionTok{c}\NormalTok{(}\FloatTok{0.025}\NormalTok{, }\FloatTok{0.5}\NormalTok{, }\FloatTok{0.975}\NormalTok{))}
\end{Highlighting}
\end{Shaded}

\begin{quote}
\textbf{Note}: the previous chunk is set to \texttt{eval=FALSE} so that
the document can be opened even if \texttt{rstan} is not installed. If
you already have it configured, you can change it to \texttt{eval=TRUE}.
\end{quote}

\section{c) 95\% credible interval for a new female
student}\label{c-95-credible-interval-for-a-new-female-student}

We seek the 95\% credible interval for the grade of a \textbf{new girl}
who did not take the exam. Under the chosen model, this corresponds to
the posterior predictive distribution of \texttt{y\_new\_girl}.

\subsection{Approximate analytical
summary}\label{approximate-analytical-summary}

As a quick approximation, using only the girls' data:

\begin{Shaded}
\begin{Highlighting}[]
\NormalTok{n\_g }\OtherTok{\textless{}{-}} \FunctionTok{length}\NormalTok{(girls)}
\NormalTok{mean\_g }\OtherTok{\textless{}{-}} \FunctionTok{mean}\NormalTok{(girls)}
\NormalTok{sd\_g }\OtherTok{\textless{}{-}} \FunctionTok{sd}\NormalTok{(girls)}

\FunctionTok{c}\NormalTok{(}\AttributeTok{n\_girls =}\NormalTok{ n\_g, }\AttributeTok{mean\_girls =}\NormalTok{ mean\_g, }\AttributeTok{sd\_girls =}\NormalTok{ sd\_g)}
\end{Highlighting}
\end{Shaded}

\begin{verbatim}
##    n_girls mean_girls   sd_girls 
##   7.000000   7.314286   1.534523
\end{verbatim}

Under the normal model with the classical noninformative prior, the
posterior predictive distribution for a new female observation is a
Student \(t\): \[
y_{\text{new},g} \mid y \sim t_{n_g-1}\left(\bar y_g,\; s_g\sqrt{1+\frac{1}{n_g}}\right).
\]

\begin{Shaded}
\begin{Highlighting}[]
\NormalTok{scale\_pred }\OtherTok{\textless{}{-}}\NormalTok{ sd\_g }\SpecialCharTok{*} \FunctionTok{sqrt}\NormalTok{(}\DecValTok{1} \SpecialCharTok{+} \DecValTok{1} \SpecialCharTok{/}\NormalTok{ n\_g)}
\NormalTok{ci\_analytical }\OtherTok{\textless{}{-}}\NormalTok{ mean\_g }\SpecialCharTok{+} \FunctionTok{qt}\NormalTok{(}\FunctionTok{c}\NormalTok{(}\FloatTok{0.025}\NormalTok{, }\FloatTok{0.975}\NormalTok{), }\AttributeTok{df =}\NormalTok{ n\_g }\SpecialCharTok{{-}} \DecValTok{1}\NormalTok{) }\SpecialCharTok{*}\NormalTok{ scale\_pred}
\NormalTok{ci\_analytical}
\end{Highlighting}
\end{Shaded}

\begin{verbatim}
## [1]  3.300189 11.328382
\end{verbatim}

This interval is a useful reference, but the main answer to the exercise
should be based on the Bayesian model fitted in Stan.

\subsection{Interval from Stan}\label{interval-from-stan}

Once the model has been fitted, the interval is obtained directly from
the posterior quantiles of \texttt{y\_new\_girl}:

\begin{Shaded}
\begin{Highlighting}[]
\NormalTok{post }\OtherTok{\textless{}{-}}\NormalTok{ rstan}\SpecialCharTok{::}\FunctionTok{extract}\NormalTok{(fit)}

\NormalTok{ci\_stan }\OtherTok{\textless{}{-}} \FunctionTok{quantile}\NormalTok{(post}\SpecialCharTok{$}\NormalTok{y\_new\_girl, }\AttributeTok{probs =} \FunctionTok{c}\NormalTok{(}\FloatTok{0.025}\NormalTok{, }\FloatTok{0.5}\NormalTok{, }\FloatTok{0.975}\NormalTok{))}
\NormalTok{ci\_stan}
\end{Highlighting}
\end{Shaded}

And we can visualize the posterior predictive distribution:

\begin{Shaded}
\begin{Highlighting}[]
\FunctionTok{hist}\NormalTok{(post}\SpecialCharTok{$}\NormalTok{y\_new\_girl, }\AttributeTok{breaks =} \DecValTok{40}\NormalTok{,}
     \AttributeTok{main =} \StringTok{"Posterior predictive distribution for a new girl"}\NormalTok{,}
     \AttributeTok{xlab =} \StringTok{"y\_new\_girl"}\NormalTok{)}
\FunctionTok{abline}\NormalTok{(}\AttributeTok{v =}\NormalTok{ ci\_stan[}\FunctionTok{c}\NormalTok{(}\DecValTok{1}\NormalTok{, }\DecValTok{3}\NormalTok{)], }\AttributeTok{lty =} \DecValTok{2}\NormalTok{)}
\end{Highlighting}
\end{Shaded}

\section{Final interpretation}\label{final-interpretation}

\subsection{On the equality between boys and
girls}\label{on-the-equality-between-boys-and-girls}

With these data, the observed difference between means is small. In the
Bayesian model fitted with Stan, the key parameter is \[
\delta = \mu_b - \mu_g.
\]

To assess whether grades are ``equal'', it is useful to examine:

\begin{itemize}
\tightlist
\item
  whether 0 lies inside the 95\% credible interval for \(\delta\),
\item
  how much posterior mass lies above or below 0,
\item
  and whether a small difference would have practical relevance.
\end{itemize}

Useful code to summarize this is:

\begin{Shaded}
\begin{Highlighting}[]
\NormalTok{ci\_delta }\OtherTok{\textless{}{-}} \FunctionTok{quantile}\NormalTok{(post}\SpecialCharTok{$}\NormalTok{delta, }\AttributeTok{probs =} \FunctionTok{c}\NormalTok{(}\FloatTok{0.025}\NormalTok{, }\FloatTok{0.5}\NormalTok{, }\FloatTok{0.975}\NormalTok{))}
\NormalTok{prob\_boys\_higher }\OtherTok{\textless{}{-}} \FunctionTok{mean}\NormalTok{(post}\SpecialCharTok{$}\NormalTok{delta }\SpecialCharTok{\textgreater{}} \DecValTok{0}\NormalTok{)}

\NormalTok{ci\_delta}
\NormalTok{prob\_boys\_higher}
\end{Highlighting}
\end{Shaded}

If the interval for \(\delta\) contains 0 comfortably, the evidence in
favor of differences between sexes will be weak.

\subsection{Compact answer for
submission}\label{compact-answer-for-submission}

\begin{enumerate}
\def\labelenumi{\arabic{enumi}.}
\item
  The model in the statement is \[
  y_i \sim \mathcal N(\beta_0 + \beta_1 x_i, \sigma),
  \] with weakly informative priors such as
  \(\beta_0 \sim \mathcal N(7,2^2)\),
  \(\beta_1 \sim \mathcal N(0,2^2)\),
  \(\sigma \sim \text{Half-Normal}(0,2)\).
\item
  That model is reasonable, but it can be improved by allowing different
  variances in each group. Therefore, we propose \[
  y_{b,k} \sim \mathcal N(\mu_b, \sigma_b), \qquad
  y_{g,j} \sim \mathcal N(\mu_g, \sigma_g),
  \] with weakly informative priors on
  \(\mu_b, \mu_g, \sigma_b, \sigma_g\).
\item
  Under this model, the 95\% credible interval for the grade of a new
  female student is obtained from the posterior predictive variable
  \texttt{y\_new\_girl} generated in Stan: \[
  \bigl[q_{0.025}(y_{\text{new},g}),\; q_{0.975}(y_{\text{new},g})\bigr].
  \] As a quick analytical approximation using the girls' data only, the
  interval is:
\end{enumerate}

\begin{Shaded}
\begin{Highlighting}[]
\NormalTok{knitr}\SpecialCharTok{::}\FunctionTok{kable}\NormalTok{(}
  \FunctionTok{data.frame}\NormalTok{(}
    \AttributeTok{method =} \StringTok{"Analytical t approximation"}\NormalTok{,}
    \AttributeTok{lower =} \FunctionTok{round}\NormalTok{(ci\_analytical[}\DecValTok{1}\NormalTok{], }\DecValTok{3}\NormalTok{),}
    \AttributeTok{upper =} \FunctionTok{round}\NormalTok{(ci\_analytical[}\DecValTok{2}\NormalTok{], }\DecValTok{3}\NormalTok{)}
\NormalTok{  ),}
  \AttributeTok{booktabs =} \ConstantTok{TRUE}\NormalTok{,}
  \AttributeTok{col.names =} \FunctionTok{c}\NormalTok{(}\StringTok{"Method"}\NormalTok{, }\StringTok{"Lower bound"}\NormalTok{, }\StringTok{"Upper bound"}\NormalTok{)}
\NormalTok{)}
\end{Highlighting}
\end{Shaded}

\begin{longtable}[]{@{}lrr@{}}
\toprule\noalign{}
Method & Lower bound & Upper bound \\
\midrule\noalign{}
\endhead
\bottomrule\noalign{}
\endlastfoot
Analytical t approximation & 3.3 & 11.328 \\
\end{longtable}

\section{Compilation instructions}\label{compilation-instructions}

\subsection{Option 1: compile the PDF from
RStudio}\label{option-1-compile-the-pdf-from-rstudio}

\begin{itemize}
\tightlist
\item
  Open \texttt{bayesian\_grades\_stan.Rmd}.
\item
  Install packages if needed:
\end{itemize}

\begin{Shaded}
\begin{Highlighting}[]
\FunctionTok{install.packages}\NormalTok{(}\FunctionTok{c}\NormalTok{(}\StringTok{"rmarkdown"}\NormalTok{, }\StringTok{"knitr"}\NormalTok{, }\StringTok{"rstan"}\NormalTok{, }\StringTok{"bayesplot"}\NormalTok{, }\StringTok{"posterior"}\NormalTok{))}
\end{Highlighting}
\end{Shaded}

\begin{itemize}
\tightlist
\item
  If you do not have LaTeX installed:
\end{itemize}

\begin{Shaded}
\begin{Highlighting}[]
\FunctionTok{install.packages}\NormalTok{(}\StringTok{"tinytex"}\NormalTok{)}
\NormalTok{tinytex}\SpecialCharTok{::}\FunctionTok{install\_tinytex}\NormalTok{()}
\end{Highlighting}
\end{Shaded}

\begin{itemize}
\tightlist
\item
  Click \textbf{Knit to PDF}.
\end{itemize}

\subsection{Option 2: from the R
console}\label{option-2-from-the-r-console}

\begin{Shaded}
\begin{Highlighting}[]
\NormalTok{rmarkdown}\SpecialCharTok{::}\FunctionTok{render}\NormalTok{(}\StringTok{"bayesian\_grades\_stan.Rmd"}\NormalTok{, }\AttributeTok{output\_format =} \StringTok{"pdf\_document"}\NormalTok{)}
\end{Highlighting}
\end{Shaded}


\end{document}
